\specialchapter{How This Book Was Made}{Colophon}{34}

\addsec{Roles}

\emph{Optimizers in the Wild} was written by Claude, an AI model made by Anthropic, in October 2026. Xiaoyu Li conceived the series \emph{The Missing Textbooks} and proposed this volume: its subject, its gap (optimization textbooks stop at convex analysis and stationary points, while practice is decided by Adam, warmup, weight decay, learning-rate schedules, Muon and the edge of stability), its two lines (the unification of optimizers as steepest descent under norms with preconditioning and moving averages, and the dynamics of training), and its position relative to existing work. Xiaoyu Li also set the direction of the series (language, design and credit) and is responsible for the book's publication.

\addsec{Process}

A written plan fixed the thesis, the fifteen chapters, the global notation, the labels of the shared appendix, and a specification for each chapter, including the editor's plan for each main theorem; several of these plans were checked numerically before writing began, and two were corrected as a result. A style guide fixed the anatomy of a chapter, the voice, a proof style that follows Halmos, the distinction between proofs, sketches, heuristic derivations, observations and folklore, and the rule that every chapter compares an exact prediction with a simulation. Separate instances of the same model then drafted the chapters and the appendix in parallel, each re-deriving its results, checking them numerically, verifying its references by web search, and writing the scripts behind its figures. A coordinating instance wrote the prologue, preface, notation, epilogue and this colophon, designed the volume's additions to the series typography, and integrated the whole.

Each chapter then received an independent referee pass from a separate instance, which re-derived the proofs, re-ran the figure scripts, checked constants and exercises numerically, and checked historical claims and references against sources wherever these could be reached. The referees corrected the errors they found directly.

\addsec{Limits}

No human has yet checked every proof. Several results are presented as heuristic derivations, marked with $\Diamond$, because no proof is known to the author or because a proof would not fit; they are checked against simulations but not proved. In the environment in which the book was written, references could be checked only through web search, not by reading the papers themselves, so bibliographic details were included only where a search confirmed them. Results from 2025 and 2026 are cited from preprints where no refereed version could be found, and are marked as such where a claim rests on them. Practices change quickly, and the recipes quoted in the book are dated. Corrections are welcome and will be recorded in the repository of \emph{The Missing Textbooks}.

\medskip
{\small\noindent Typeset with pdf\LaTeX{} in \texttt{missingbook}, the series class built on \textsc{koma}-Script, with the volume style \texttt{wild}; Palatino (newpx) for text and mathematics, biblatex and biber for references, Ti\emph{k}Z and pgfplots for figures, whose data were computed with Python, NumPy, SciPy and PyTorch. The emblem on the title page shows one gradient and its three steepest-descent directions, the points of the $\ell_1$, $\ell_2$ and $\ell_\infty$ unit balls that make the largest inner product with it; the wave at the foot of each opening page rises from chapter to chapter.\par}
